\scanpage{17}

\deleted{Supposons \unsure{que tous les corps résiduels de}}

Notons qu'il y a dans $C$ un $S$-automorphisme canonique de symétrie, car $C$ est un schéma abélien sur $S$. Cet automorphisme \inserted{est d'ordre $2$, il} opère sur les $\mathcal E_i$ [de façon compatible avec les structures, multiplications \inserted{et les opérations de $\mathcal L$ sur $\mathcal E_\infty$}] et induit sur $\mathcal L$ le chgt de signe.
\deleted{\unsure{Il en résulte aussitôt qu'il opère trivialement dans $\mathcal E_2$ [car il y opère trivialement]}}
\deleted{Si les caractéristiques résiduelles de $S$ sont $\neq2$, on en conclut}

Supposons les car. résiduelles de $S\neq2$. Alors
\deleted{Il existe un supplémentaire unique de $\mathcal E_1$ dans $\mathcal E_2$ \unsure{stable qui sera le lieu $s$ opère}}
$s$ opère trivialement dans $\mathcal E_2$ [c'est ainsi \unsure{sans|sous} la condition qu'on vient de faire, on le verra], et il existe un supplémentaire unique \inserted{$L_3$} de $\mathcal E_2$ dans $\mathcal E_3$ où $s$ opère par symétrie, \deleted{\unsure{le}} qui est aussi égal à $\mathcal L\cdot \mathcal E_2$, on le note $L_3$.
\margin{N.B. la \unsure{considération} de $C\longrightarrow P(\mathcal E_2)$ et de la \unsure{différentielle} de ce morphisme nous donne (sans restriction de car) un morphisme
\diagram{\begin{tikzcd}\unsure{\mathcal L|L}^{\otimes3}\arrow[r,"\unsure{\partial|\vartheta}"]&\mathcal E_3\end{tikzcd}}
\unsure{compatible avec tt changement de base $S$}.}

\deleted{Alors $L_3^{\otimes2}$ est un supplémentaire de $\mathcal E_5$ dans $\mathcal E_6$, soit $L_6$. A condition}
On peut donc normaliser encore plus les bases admissibles de $\mathcal E_3$, en exigeant que $x$ soit une base de $\unsure{\mathcal E_3|L_3}$ ; donc le groupe des changements de base admis est celui de la forme
\diagram{\[\left(\begin{array}{ccc}1&\alpha&0\\0&\lambda^2&0\\0&0&\lambda^3\end{array}\right)\qquad\text{i.e.}\qquad\left(\begin{array}{ccc}1&\alpha&0\\0&u&0\\0&0&v\end{array}\right)\]}
avec $\deleted{\unsure{\lambda}}u^3-v^2=0$.
Ce n'est autre que le \uline{groupe affine à $1$ variable}.

Quelle est la normalisation correspondante de l'équation de $C$ ?
\mspage{16}
